\documentclass[../../main.tex]{subfiles}

\begin{document}

\section{Lebesgue 外测度}\label{section:ch02-外测度}

令 $I$ 为实数的非空区间. 若 $I$ 是无界的, 定义它的长度 $\ell(I)$ 为 $\infty$, 否则定义它的长度为端点的差. 对于一个实数集合 $A$, 考虑覆盖 $A$ 的非空开有界区间的可数集族 $\{I_{k}\}_{k=1}^{\infty}$, 即使得 $A\subseteq\bigcup_{k=1}^{\infty}I_{k}$ 的族. 对每个这样的族, 考虑该族的区间的长度之和. 由于长度是正数, 每个和是唯一定义的且与项的顺序无关. 我们定义 $A$ 的外测度 $m^{*}(A)$ 为所有这样的和的下确界, 即
\begin{align}
m^{*}(A)=\inf\left\{\sum_{k=1}^{\infty}\ell(I_{k})\ \middle|\ A\subseteq\bigcup_{k=1}^{\infty}I_{k}\right\}.\label{align:ch02-外测度定义}
\end{align}

从外测度的定义立即得出 $m^{*}(\varnothing)=0$. 此外, 由于集合 $B$ 的任何覆盖也是 $B$ 的任何子集的覆盖, 所以外测度在以下意义下是单调的:
\begin{align*}
\text{若 }A\subseteq B,\quad\text{则 }m^{*}(A)\leqslant m^{*}(B).
\end{align*}

\begin{example}\label{example:ch02-可数集零外测度}
可数集具有零外测度. 事实上, 令 $C$ 为可数集, 列举为 $C=\{c_{k}\}_{k=1}^{\infty}$. 令 $\varepsilon>0$. 对每个自然数 $k$, 定义 $I_{k}=(c_{k}-\varepsilon/2^{k+1},c_{k}+\varepsilon/2^{k+1})$. 开区间的可数族 $\{I_{k}\}_{k=1}^{\infty}$ 覆盖 $C$. 因此
\begin{align*}
0\leqslant m^{*}(C)\leqslant\sum_{k=1}^{\infty}\ell(I_{k})=\sum_{k=1}^{\infty}\varepsilon/2^{k}=\varepsilon.
\end{align*}
不等式对每个 $\varepsilon>0$ 成立. 因此 $m^{*}(C)=0$.
\end{example}

\begin{proposition}[区间的外测度是它的长度]\label{proposition:ch02-区间外测度}
\end{proposition}

\begin{proof}
我们从闭有界区间 $[a,b]$ 的情形开始. 令 $\varepsilon>0$. 由于开区间 $(a-\varepsilon,b+\varepsilon)$ 包含 $[a,b]$, 我们有 $m^{*}([a,b])\leqslant\ell((a-\varepsilon,b+\varepsilon))=b-a+2\varepsilon$. 这对任何 $\varepsilon>0$ 成立. 因此 $m^{*}([a,b])\leqslant b-a$. 接下来要证明 $m^{*}([a,b])\geqslant b-a$. 而这等价于证明: 若 $\{I_{k}\}_{k=1}^{\infty}$ 是任何覆盖 $[a,b]$ 的可数开有界区间族, 则
\begin{align}
\sum_{k=1}^{\infty}\ell(I_{k})\geqslant b-a.\label{align:ch02-覆盖下界}
\end{align}

根据 Heine-Borel 定理（\cref{theorem:ch01-HeineBorel}）, 任何覆盖 $[a,b]$ 的开区间族有一个覆盖 $[a,b]$ 的有限子族. 选取自然数 $n$ 使得 $\{I_{k}\}_{k=1}^{n}$ 覆盖 $[a,b]$. 我们将证明
\begin{align}
\sum_{k=1}^{n}\ell(I_{k})\geqslant b-a.\label{align:ch02-有限覆盖下界}
\end{align}
从而\eqref{align:ch02-覆盖下界}成立. 由于 $a$ 属于 $\bigcup_{k=1}^{n}I_{k}$, 这些 $I_{k}$ 中必有一个包含 $a$. 选取这样的一个区间且记为 $(a_{1},b_{1})$. 我们有 $a_{1}<a<b_{1}$. 若 $b_{1}\geqslant b$, 不等式\eqref{align:ch02-有限覆盖下界}得证, 这是因为
\begin{align*}
\sum_{k=1}^{n}\ell(I_{k})\geqslant b_{1}-a_{1}>b-a.
\end{align*}
否则, $b_{1}\in[a,b]$, 且由于 $b_{1}\notin(a_{1},b_{1})$, 族 $\{I_{k}\}_{k=1}^{n}$ 中存在一个区间, 记为 $(a_{2},b_{2})$ 以区别于 $(a_{1},b_{1})$, 使得 $b_{1}\in(a_{2},b_{2})$, 即 $a_{2}<b_{1}<b_{2}$. 若 $b_{2}\geqslant b$, 不等式\eqref{align:ch02-有限覆盖下界}得证, 这是因为
\begin{align*}
\sum_{k=1}^{n}\ell(I_{k})\geqslant(b_{1}-a_{1})+(b_{2}-a_{2})=b_{2}-(a_{2}-b_{1})-a_{1}>b_{2}-a_{1}>b-a.
\end{align*}
我们继续这一选取程序直至它终止, 而它必须终止, 因为族 $\{I_{k}\}_{k=1}^{n}$ 中仅有 $n$ 个区间. 因此我们得到 $\{I_{k}\}_{k=1}^{n}$ 的一个子族 $\{(a_{k},b_{k})\}_{k=1}^{N}$ 使得
\begin{align*}
a_{1}<a,\qquad\text{而对 }1\leqslant k\leqslant N-1,\quad a_{k+1}<b_{k},
\end{align*}
且由于选取过程终止, $b_{N}>b$. 因此
\begin{align*}
\sum_{k=1}^{n}\ell(I_{k})&\geqslant(b_{N}-a_{N})+(b_{N-1}-a_{N-1})+\cdots+(b_{1}-a_{1})\\
&=b_{N}-(a_{N}-b_{N-1})-\cdots-(a_{2}-b_{1})-a_{1}\\
&>b_{N}-a_{1}>b-a.
\end{align*}
因而不等式\eqref{align:ch02-有限覆盖下界}成立.

若 $I$ 是任意有界区间, 则给定 $\varepsilon>0$, 存在两个闭有界区间 $J_{1}$ 和 $J_{2}$ 使得
\begin{align*}
J_{1}\subseteq I\subseteq J_{2},
\end{align*}
而
\begin{align*}
\ell(I)-\varepsilon<\ell(J_{1})\quad\text{且}\quad\ell(J_{2})<\ell(I)+\varepsilon.
\end{align*}
根据对闭有界区间的外测度与长度的相等性, 以及外测度的单调性, 有
\begin{align*}
\ell(I)-\varepsilon<\ell(J_{1})=m^{*}(J_{1})\leqslant m^{*}(I)\leqslant m^{*}(J_{2})=\ell(J_{2})<\ell(I)+\varepsilon.
\end{align*}
这对每个 $\varepsilon>0$ 成立. 因此 $\ell(I)=m^{*}(I)$.

若 $I$ 是无界区间, 则对每个自然数 $n$, 存在区间 $J\subseteq I$ 满足 $\ell(J)=n$. 因此 $m^{*}(I)\geqslant m^{*}(J)=\ell(J)=n$. 这对每个自然数 $n$ 成立, 因此 $m^{*}(I)=\infty$.
\end{proof}

\begin{proposition}[外测度平移不变]\label{proposition:ch02-外测度平移}
外测度是平移不变的, 即对任意集合 $A$ 与数 $y$,
\begin{align*}
m^{*}(A+y)=m^{*}(A).
\end{align*}
\end{proposition}

\begin{proof}
观察到若 $\{I_{k}\}_{k=1}^{\infty}$ 是任意可数集族, 则 $\{I_{k}\}_{k=1}^{\infty}$ 覆盖 $A$ 当且仅当 $\{I_{k}+y\}_{k=1}^{\infty}$ 覆盖 $A+y$. 此外, 若每个 $I_{k}$ 是一个开区间, 则每个 $I_{k}+y$ 是一个相同长度的开区间, 因而
\begin{align*}
\sum_{k=1}^{\infty}\ell(I_{k})=\sum_{k=1}^{\infty}\ell(I_{k}+y).
\end{align*}
结论从这两个观察可以得到.
\end{proof}

\begin{proposition}[外测度可数次可加]\label{proposition:ch02-外测度次可加}
外测度是可数次可加的, 即若 $\{E_{k}\}_{k=1}^{\infty}$ 是任意可数集族, 互不相交或相交, 则
\begin{align*}
m^{*}\left(\bigcup_{k=1}^{\infty}E_{k}\right)\leqslant\sum_{k=1}^{\infty}m^{*}(E_{k}).
\end{align*}
\end{proposition}

\begin{proof}
若这些 $E_{k}$ 中的一个有无穷的外测度, 则不等式平凡地成立. 我们因此假定每个 $E_{k}$ 有有限的外测度. 令 $\varepsilon>0$. 对每个自然数 $k$, 存在开有界区间的可数族 $\{I_{k,i}\}_{i=1}^{\infty}$ 使得
\begin{align*}
E_{k}\subseteq\bigcup_{i=1}^{\infty}I_{k,i}\quad\text{且}\quad\sum_{i=1}^{\infty}\ell(I_{k,i})<m^{*}(E_{k})+\varepsilon/2^{k}.
\end{align*}
现在 $\{I_{k,i}\}_{1\leqslant k,i<\infty}$ 是一个覆盖 $\bigcup_{k=1}^{\infty}E_{k}$ 的开有界区间的可数族: 由于该族是可数族组成的可数族, 它是可数的. 因此, 根据外测度的定义,
\begin{align*}
m^{*}\left(\bigcup_{k=1}^{\infty}E_{k}\right)\leqslant\sum_{1\leqslant k,i<\infty}\ell(I_{k,i})=\sum_{k=1}^{\infty}\left[\sum_{i=1}^{\infty}\ell(I_{k,i})\right]<\sum_{k=1}^{\infty}\left[m^{*}(E_{k})+\varepsilon/2^{k}\right]=\left[\sum_{k=1}^{\infty}m^{*}(E_{k})\right]+\varepsilon.
\end{align*}
由于这对每个 $\varepsilon>0$ 成立, 它对 $\varepsilon=0$ 也成立. 证明完毕.
\end{proof}

若 $\{E_{k}\}_{k=1}^{n}$ 是任何有限集族, 互不相交或相交, 则
\begin{align*}
m^{*}\left(\bigcup_{k=1}^{n}E_{k}\right)\leqslant\sum_{k=1}^{n}m^{*}(E_{k}).
\end{align*}
通过对 $k>n$ 设 $E_{k}=\varnothing$, 有限次可加性从可数次可加性得到.

\begin{exercise}\label{exercise:ch02-05}
用外测度的性质证明区间 $[0,1]$ 不是可数的.
\end{exercise}

\begin{exercise}\label{exercise:ch02-06}
令 $A$ 为区间 $[0,1]$ 的无理数集. 证明 $m^{*}(A)=1$.
\end{exercise}

\begin{exercise}\label{exercise:ch02-07}
实数集称为 $G_{\delta}$ 集, 若它是可数个开集族的交. 证明对任何有界集 $E$, 存在 $G_{\delta}$ 集 $G$ 使得 $E\subseteq G$ 且 $m^{*}(G)=m^{*}(E)$.
\end{exercise}

\begin{exercise}\label{exercise:ch02-08}
令 $B$ 为区间 $[0,1]$ 内的无理数集, 令 $\{I_{k}\}_{k=1}^{n}$ 为覆盖 $B$ 的有限开区间族. 证明 $\sum_{k=1}^{n}m^{*}(I_{k})\geqslant 1$.
\end{exercise}

\begin{exercise}\label{exercise:ch02-09}
证明: 若 $m^{*}(A)=0$, 则 $m^{*}(A\cup B)=m^{*}(B)$.
\end{exercise}

\begin{exercise}\label{exercise:ch02-10}
令 $A$ 和 $B$ 为有界集使得存在 $\alpha>0$, 对所有 $a\in A$, $b\in B$, $|a-b|\geqslant\alpha$. 证明 $m^{*}(A\cup B)=m^{*}(A)+m^{*}(B)$.
\end{exercise}

\end{document}